\section{Conclusion}\label{section: conclusion}

We have shown that the star-height hierarchy of Kleene algebra with complement is infinite \kl{wrt} languages, already for \kl{terms} over four variables (\Cref{theorem: main}), and also over the finite relational models $\RSUB$ and over finite relational models.
The separating invariant is arithmetic: the degree of the irreducible factors of the polynomial that describes the periods of a \kl{term} in parameterized \kl{chain models}.
To propagate this invariant through nested stars, we have shown that the weights of walks in graphs labeled by ultimately periodic \kl{parametric Presburger families} are again parametric Presburger (\Cref{lemma: closure}), which may be of independent interest.

We conclude with some open questions.
\begin{itemize}
    \item \emph{Fewer variables.} Does the hierarchy remain infinite for \kl{terms} with fewer than four variables, in particular with a single variable?
    \item \emph{Restricted model classes.} Our counter-models have unbounded treewidth. Is the hierarchy still infinite over structures of bounded treewidth, \eg, over series-parallel structures?
    \item \emph{Standard interpretation.} Under the standard interpretation, every regular language has generalized star height at most three~\cite{openaiGeneralizedStarHeight2026}, and it remains open whether height two is ever needed~\cite{pinResultsGeneralizedStarheight1992}. Which intermediate semantics, between arbitrary languages and letters, still yield an infinite hierarchy?
    \item \emph{Decidability.} Is it decidable, given a \kl{term} $\term$ and $h \in \nat$, whether $\term$ is \kl{language-equivalent} to a \kl{term} of \kl{star height} at most $h$? Since \kl{language-equivalence} itself is $\Pi^{0}_{1}$-complete~\cite{nakamuraFiniteRelationalSemantics2025}, the exact complexity of this problem is an interesting question; already the case $h = 0$ (definability by a star-free \kl{term}) is open to our knowledge.
\end{itemize}
